\begin{proofbox}
	\textbf{\textcolor{CProof}{思路提示}}

	本结论可以通过两条已知定理直接贯通：正弦定理将边与正弦值挂钩，大边对大角将边与角挂钩。因此只需证明三者中任一对等价，再串联即可。

	\medskip
	\textbf{\textcolor{CProof}{正式推导}}
	\begin{description}[style=nextline, leftmargin=2.8em, labelsep=0.5em, itemsep=0.55em]
		\item[\textbf{步骤一（正弦定理转化）：}] 证明 $a>b\iff\sin A>\sin B$。
			\[
				\begin{aligned}
					\because \quad & \frac{a}{\sin A} = 2R,\\
					& \frac{b}{\sin B} = 2R,\\
					\therefore \quad & a = 2R\sin A,\quad b = 2R\sin B.
			\end{aligned}
			\]
			\par\textbf{理由：} 由 $2R>0$ 知 $a>b\iff\sin A>\sin B$，正弦定理建立了边长与正弦值的线性关系，消去 $2R$ 直接比较。

		\item[\textbf{步骤二（角度正相关）：}] 证明 $A>B\iff\sin A>\sin B$。

			先证 $A>B\Rightarrow\sin A>\sin B$。分两类：
			\begin{itemize}
			\item 若 $A,B\in(0,\frac{\pi}{2}]$，正弦函数在该区间递增，故 $\sin A>\sin B$。
			\item 若 $A\in(\frac{\pi}{2},\pi)$，$B\in(0,\frac{\pi}{2})$，由 $A+B<\pi$ 得 $B<\pi-A<\frac{\pi}{2}$。而 $\sin A=\sin(\pi-A)$ 且 $\sin$ 在 $(0,\frac{\pi}{2})$ 上增，因此 $\sin(\pi-A)>\sin B$，即 $\sin A>\sin B$。
		\end{itemize}
			再证反向：若 $\sin A>\sin B$，假设 $A<B$，则由正向结论（交换 $A,B$）可得 $\sin B>\sin A$，矛盾；若 $A=B$，则 $\sin A=\sin B$，矛盾。故 $A>B$。
			\[
				A>B\iff\sin A>\sin B.
			\]
			\par\textbf{理由：} 利用正弦函数单调区间及三角形内角限制，分类讨论并反证。

		\item[\textbf{步骤三（等价链合并）：}] 将两个等价关系串联。
			\[
				a>b\iff\sin A>\sin B\iff A>B.
			\]
			\par\textbf{理由：} 等价的传递性。
	\end{description}

	\medskip
	\textbf{\textcolor{CProof}{结论回扣}}

	\textbf{由此建立了边、角、正弦三者的无缝等价链，在三角形范围内可互相代换。}
\end{proofbox}
